% Requires: amsmath (standard in most thesis templates).
% Notation: adjust \Delta (cell size) and symbol names to match the rest of
% your thesis.

\section{Release Area Estimation from Mapped Avalanche Outlines}
\label{sec:release_estimation}

Observed avalanche outlines usually document the full extent of an event, but
numerical simulations require the release area at the top of the path, which is
often not mapped separately. This section describes an automated method that
derives a plausible release area from the outline polygon and a digital
elevation model (DEM). The method proceeds in two steps: first the crown line
--- the upper fracture edge of the slab --- is detected from the terrain, then
a release area covering a prescribed fraction of the outline area (typically
20--30\,\%) is filled in below it.

Throughout, the outline is a mask $M$ of grid cells, and $z(c)$ is the elevation
of a cell $c$ of the DEM with cell size $\Delta$.

\subsection{Crown Line Detection}
\label{sec:crown_detection}

The crown line is the part of the outline boundary where the terrain enters the
polygon from above. The method identifies it with flow directions, not with
elevation contours, so it also works when the crown line crosses contours at an
angle.

\paragraph{Flow routing.}
A water droplet on each cell is assumed to move to the steepest downslope of
the eight neighbouring cells (the \emph{D8} direction),
\begin{equation}
  \mathrm{flow}(c)
  = \arg\max_{n \,\in\, N_8(c)} \frac{z(c) - z(n)}{\lVert n - c \rVert},
  \label{eq:d8}
\end{equation}
where $N_8(c)$ is the 8-neighbourhood of $c$ and $\lVert n - c \rVert$ is
$\Delta$ for orthogonal and $\sqrt{2}\,\Delta$ for diagonal neighbours. Before
routing, closed depressions in the DEM are filled so that every cell has a
downslope path out of the domain; this is the same sink filling used in
standard hydrological terrain analysis.

\paragraph{Inflow boundary.}
A cell of the outline is a \emph{crown candidate} if at least one cell outside
the outline drains into it,
\begin{equation}
  c \in M \text{ is a crown candidate}
  \quad\Longleftrightarrow\quad
  \exists\, n \notin M : \; \mathrm{flow}(n) = c .
  \label{eq:crown}
\end{equation}
These cells form the upstream boundary of the polygon: the crown line at the
top, but possibly also patches on the flanks where side gullies cross the
outline, or on the valley floor where the opposite slope drains into the
avalanche track. The candidates are grouped into connected components $K$, and
each component is tested for steepness: a slab avalanche releases on steep
terrain, while inflow from an opposite slope ends on the gentle valley floor or
in the runout. With the slope angle $\theta(c) = \arctan \lVert \nabla z(c)
\rVert$ from central differences, components with
\begin{equation}
  \frac{1}{|K|} \sum_{c \in K} \theta(c) < \theta_{\min}
  \label{eq:crown_slope}
\end{equation}
are discarded, where $\theta_{\min}$ is the lower bound of the release slope
window (28$^\circ$ by default). On two mapped outlines, the crown components
had mean slopes of 35--71$^\circ$ and the valley-floor components 12--22$^\circ$.
Unlike a threshold on absolute elevation, this criterion keeps secondary crown
lines of lower-lying lobes. Optionally, components whose mean elevation
lies below a given share of the outline's elevation range can be discarded as
well (disabled by default). Components smaller than three cells are always
dropped. If no component survives, the largest one is kept as a fallback.

\paragraph{Particle-based variant.}
Instead of the D8 directions of Eq.~\eqref{eq:d8}, the drainage into the
outline can also be measured with the avalanche model itself. A separate simulation is run
on a copy of the DEM in which the outline is masked out (set to no-data): mass
points are released on all terrain outside the outline with the plain sliding
model, i.e.\ without curvature and without particle interactions, and they
stop as soon as the terrain sampling at their position touches the masked
outline. Each stopped point is attributed to the outline cell in the direction
it was moving, and outline cells entered by at least a small number of points
(two, at four points per release cell) become crown candidates, which are then
filtered as above. Because the points carry inertia, this variant also detects
terrain that reaches the outline only with momentum, for example across a
counter-slope shoulder; it is, however, roughly as expensive as one additional
avalanche simulation.

\paragraph{Gradient-step variant.}
A purely geometric alternative avoids routing altogether. From the centre
$\mathbf{x}_c$ of every outline cell, a step of one cell length $\Delta$ is
taken uphill, along the normalized elevation gradient (central differences),
\begin{equation}
  \mathbf{p}_c = \mathbf{x}_c + \Delta \, \frac{\nabla z(c)}{\lVert \nabla z(c) \rVert} .
  \label{eq:gradient_step}
\end{equation}
If $\mathbf{p}_c$ still lies in an outline cell, $c$ is interior. Otherwise the
distance $d_c$ from $\mathbf{p}_c$ to the nearest outline cell centre is
measured, and $c$ is a crown candidate if $d_c > \delta\,\Delta$ with a
threshold $\delta \in [0, 1)$ (default $\delta = 0.75$). A step straight
out of the upstream edge lands $d_c = \Delta$ away; every step that leaves the
outline lands at least $\Delta/2$ away. The threshold is needed because
avalanches spread laterally on their way down: on the flanks of a widening
outline, the uphill step leaves the polygon at a grazing angle and ends up only
slightly outside it, although these cells are not part of the crown line.
Isolated stair-step cells of a rasterized flank exceed the threshold, but they
form tiny components that the component filter above removes. The candidates
are filtered exactly as in the flow-routing method. The variant is the cheapest
of the three, but it uses only the local slope direction, so terrain that
converges onto the outline from further away is not considered.

\subsection{Geodesic Fill}
\label{sec:geodesic_fill}

From the crown line, the release area is grown downhill until it covers a
target fraction $f$ of the outline area (here $f = 0.25$). For every outline
cell $c$, the geodesic distance $d(c)$ to the crown line is computed: the
length of the shortest path from the crown line to $c$ that stays inside the
outline, counting $\Delta$ per orthogonal and $\sqrt{2}\,\Delta$ per diagonal
step. Only cells that could release at all are eligible: the potential
release areas $P \subseteq M$ are the outline cells that pass the
terrain classification also used for automatic release areas, i.e.\ a slope
between $\theta_{\min}$ and $\theta_{\max}$, an elevation above a minimum
release elevation and a surface roughness below a threshold. The release area
$R$ is then defined as the eligible cells with the smallest distances,
\begin{equation}
  R = \text{the } \min\!\big(\operatorname{round}(f \, |M|),\, |P|\big)
      \text{ cells of } P \text{ with the smallest } d(c),
  \label{eq:fill}
\end{equation}
where $|M|$ is the number of outline cells. All cells of $R$ are assigned a
uniform slab thickness. The distance is still measured through the whole
outline, so a rock band or a flat bench directly below the crown is skipped
without cutting off the release terrain behind it. Because the fill front
advances at equal distance from the crown line, it remains \emph{parallel to
the crown line}: a band of uniform width below the fracture edge, independent
of the contour lines. This is the key difference from selecting cells above an
elevation threshold, which would fail whenever the crown line does not follow
the contours. If the potential release areas are smaller than the target, all
reachable ones are selected and the release fraction falls short of $f$.

\subsection{Plausibility Checks and Limitations}
\label{sec:release_checks}

The method reports the mean slope of the detected crown line, computed from the
steepest descent on the raw DEM. If it falls outside the typical release
window of about 28--60$^\circ$, the outline most likely documents only the
deposition zone and the estimated release area should not be trusted. Outline
cells that cannot be reached from the crown line (disconnected polygon parts)
are never selected. The fraction $f$ is the main tuning parameter: it should
grow with event size, and 20--30\,\% is a sensible default for mid-sized to
large avalanches. The slope filter of Section~\ref{sec:crown_detection} does
not remove inflow on steep lower flanks; for such cases the optional elevation
filter remains available.
